\section{Storage and Data Integrity}\label{sec:storage}

\subsection{Live pool health}

Five ZFS pools were inspected live (E-02, E-03, E-04). Four are healthy; one carries a
faulted device.

{\footnotesize\sansfont
\begin{tabularx}{\linewidth}{@{}p{2.0cm}p{1.5cm}p{1.2cm}p{3.0cm}X@{}}
\toprule
\rh{Pool} & \rh{Host} & \rh{Size} & \rh{Geometry} & \rh{State} \\
\midrule
datastore & Randy & 36.7\,T & 4$\times$ RAIDZ2 (three 6-wide + one 4-wide Seagate) & ONLINE, 0 errors \\
bulk & Randy & 80.0\,T & 3$\times$ RAIDZ2 (8+8+6-wide) & DEGRADED - \code{mpathv} FAULTED \\
workspace & QuarkyLab & 10.9\,T & 6-wide RAIDZ1 + 1 hot spare (PERC H330 JBOD) & ONLINE, 0 errors \\
tank & Jarvis & 9.09\,T & 5-wide RAIDZ1 (HBA330 IT) & ONLINE, 0 errors \\
scratch & Jarvis & 186\,G & single SSD & ONLINE, 0 errors \\
\end{tabularx}}

The live \code{datastore} layout confirms the disputed documentation fact D-04: the fourth
vdev (\code{raidz2-3}) is four Seagate drives by-WWN, in-pool, with no unallocated spares.
The \code{bulk} geometry confirms D-05: three vdevs totalling 80\,T, superseding the stale
"2$\times$8-wide / 58.2\,T" still in topology and \code{Proxmox Cluster.md}.

\subsection{The faulted bulk drive}\label{sub:mpathv}

\sevmed\; \textbf{QL-R01 - reduced parity, corrected impact.}
\code{zpool status bulk} (E-04) shows device \code{mpathv} in vdev \code{raidz2-2} as
\code{FAULTED} with \code{3} read / \code{14} write errors and "too many errors". The device
is a SEAGATE ST4000NM0023 (4\,TB SAS), WWN \code{5000c500631a54fb}; both SAS paths are dead,
so this is a failed physical drive, not a multipath flap. It corresponds to the DS4246 slot-15
drive that failed 2026-07-21 (incident NF-INC-2026-07-22). Follow-up inspection confirmed the
drive is now fully offline: it has no \code{/dev/mapper} entry at all, and ZFS is holding it
only by cached label.

Key mitigating facts, all live-verified: the pool reports \code{errors: No known data errors};
RAIDZ2 tolerates one drive loss per vdev, so \code{raidz2-2} still has one parity drive of
protection remaining; the last scrub (Sun Jul 19) completed with zero errors; and two DS4246
bays are free for an immediate \code{zpool replace}.

\begin{callout}[title={Correction: this is loss of redundancy, not imminent data loss}]
An earlier draft of this review rated the drive urgent on the basis that \code{bulk} held
researcher \code{fernanda}'s DUNE data with no second copy. \textbf{Live inspection (E-29)
contradicts that}, and the original framing came from stale documentation. \code{zfs list -r
bulk} shows \code{bulk/fernanda}, \code{bulk/media}, \code{bulk/misc} and \code{bulk/archive}
are all 205 to 222\,KB placeholders. The only real data is \code{bulk/backups} at 57.3\,GB,
which holds the Ares restic repository and a \code{jarvis-netframe} backup - both themselves
secondary copies whose sources still exist. The DUNE research data has \textbf{not} been
migrated onto \code{bulk}.

Replacing \code{mpathv} is therefore prudent maintenance to restore full double parity, not a
race against data loss. Two conditions keep it on the list rather than closing it: the
exposure becomes real the moment research data migrates onto \code{bulk}, so the second copy
(QL-R02) must exist \emph{before} that migration; and the Ares restic backup currently targets
this degraded pool. Recommended action QL-REC01, reclassified to next maintenance window and
pending operator hardware.
\end{callout}

\subsection{Controllers and multipath topology}

Randy uses two distinct controllers by design: an AVAGO 3108 MegaRAID for the boot mirror
(relocated 2026-07-01 after its original slot died - do not reuse the dead slot) and an
LSI 9207-8e in IT mode for the DS4246 data path. The DS4246 is dual-SAS cabled with
\code{multipath}, so every shelf drive presents as an \code{mpathX} device. The single
enclosure is itself a single point of failure for \code{bulk} (H-09, QL-R05 context):
loss of the shelf or its cabling takes the whole pool offline.

\subsection{Redundancy posture and aged media}

The two GPU-node pools are RAIDZ1 (QL-R21): \code{workspace} tolerates one loss but carries a
hot spare; \code{tank} tolerates one loss with no spare; \code{scratch} is a single SSD with
no redundancy (scratch by intent). Several disks are aged SAS. RAIDZ1 on old disks carries
resilver-window risk - a second fault during rebuild loses the pool. This is an accepted risk
today given clean scrubs and low utilisation, but future rebuilds should prefer RAIDZ2 and
keep spares on hand.

\subsection{Data integrity and capacity}

Scrubs are scheduled and recent (datastore and bulk both scrubbed Jul 19, zero repaired).
Scrutiny monitors roughly 79 drives across Randy, QuarkyLab, and Jarvis. The alerting
question paired with this pool is now settled: the gap was not Scrutiny or delivery but the
absence of any device-level ZFS metric, since closed (QL-R11, Section~\ref{sec:monitoring}). Utilisation is low across the board -
\code{datastore} 176\,G of 36.7\,T, \code{bulk} 80.7\,G of 80\,T, \code{workspace} 144\,G of
10.9\,T - so there is no near-term capacity pressure; the risk is redundancy and backup, not
space.

\subsection{Recoverability by dataset}

Guest images and configs are recoverable from PBS within roughly 24\,h RPO (Section~\ref{sec:backup}).
The exceptions are the two datasets with no independent copy: \code{bulk} (which today holds
only \code{bulk/backups}, research data having not yet migrated) and pve1's CT103/CT104. Ares' own \code{restic} backup, notably, targets
\code{sftp:randy:/mnt/bulk/backups/ares} - it lands on the degraded \code{bulk} pool, so the
workstation backup shares the same at-risk failure domain and should be repointed once a
second target exists.
